% Copyright (c) 2026, Ali-Eimaan. All rights reserved.
% SPDX-License-Identifier: BSD-3-Clause

% Shared preamble for every document in docs/derivations/.
% Included via % Copyright (c) 2026, Ali-Eimaan. All rights reserved.
% SPDX-License-Identifier: BSD-3-Clause

% Shared preamble for every document in docs/derivations/.
% Included via % Copyright (c) 2026, Ali-Eimaan. All rights reserved.
% SPDX-License-Identifier: BSD-3-Clause

% Shared preamble for every document in docs/derivations/.
% Included via % Copyright (c) 2026, Ali-Eimaan. All rights reserved.
% SPDX-License-Identifier: BSD-3-Clause

% Shared preamble for every document in docs/derivations/.
% Included via \input{preamble}. Define macros here only -- never in an individual file,
% or the three documents will drift into inconsistent notation.
%
% Notation single-source-of-truth: docs/README_math.md section 1. Every macro below
% corresponds one-to-one to a row of that table; do not add a symbol there without a
% macro here (or vice versa).

\usepackage[margin=1in]{geometry}
\usepackage{amsmath, amssymb, amsthm}
\usepackage{mathtools}
\usepackage{bm}
\usepackage{algorithm}
\usepackage{algpseudocode}
\usepackage{hyperref}
\usepackage{cleveref}

% --- theorem environments ---------------------------------------------------------
\theoremstyle{plain}
\newtheorem{theorem}{Theorem}
\newtheorem{lemma}[theorem]{Lemma}
\newtheorem{proposition}[theorem]{Proposition}
\newtheorem{corollary}[theorem]{Corollary}
\theoremstyle{definition}
\newtheorem{definition}[theorem]{Definition}
\newtheorem{assumption}{Assumption}
\theoremstyle{remark}
\newtheorem{remark}[theorem]{Remark}

% --- notation ---------------------------------------------------------------------
% One macro per symbol in docs/README_math.md section 1.

\newcommand{\R}{\mathbb{R}}
\newcommand{\Nbr}{\mathcal{N}}                   % open neighborhood (subscript written as \Nbr_i)
\newcommand{\Nbrc}{\bar{\mathcal{N}}}            % closed neighborhood (subscript written as \Nbrc_i)
\newcommand{\ycopy}[2]{y_{#1}^{#2}}             % local copy
\newcommand{\zcon}[1]{z^{#1}}                   % consensus variable
\newcommand{\dual}[2]{\lambda_{#1}^{#2}}        % scaled dual (code `lam`)
\newcommand{\udual}[2]{\nu_{#1}^{#2}}           % unscaled dual = rho * scaled
\newcommand{\Lagr}{\mathcal{L}}
\newcommand{\Lrho}{\Lagr_\rho}
\newcommand{\norm}[1]{\left\lVert #1 \right\rVert}
\newcommand{\prox}{\operatorname{prox}}
\newcommand{\grad}{\nabla}
\newcommand{\diag}{\operatorname{diag}}
\newcommand{\tr}{\operatorname{tr}}
\newcommand{\Eset}{\mathcal{E}}                 % formation-edge set
\newcommand{\Gset}{\mathcal{G}}                 % communication graph
\newcommand{\Fiedler}{\lambda_2(L)}             % algebraic connectivity
\newcommand{\lmax}{\lambda_{\max}}
\newcommand{\lmin}{\lambda_{\min}}
\DeclareMathOperator*{\argmin}{arg\,min}
\DeclareMathOperator*{\argmax}{arg\,max}
. Define macros here only -- never in an individual file,
% or the three documents will drift into inconsistent notation.
%
% Notation single-source-of-truth: docs/README_math.md section 1. Every macro below
% corresponds one-to-one to a row of that table; do not add a symbol there without a
% macro here (or vice versa).

\usepackage[margin=1in]{geometry}
\usepackage{amsmath, amssymb, amsthm}
\usepackage{mathtools}
\usepackage{bm}
\usepackage{algorithm}
\usepackage{algpseudocode}
\usepackage{hyperref}
\usepackage{cleveref}

% --- theorem environments ---------------------------------------------------------
\theoremstyle{plain}
\newtheorem{theorem}{Theorem}
\newtheorem{lemma}[theorem]{Lemma}
\newtheorem{proposition}[theorem]{Proposition}
\newtheorem{corollary}[theorem]{Corollary}
\theoremstyle{definition}
\newtheorem{definition}[theorem]{Definition}
\newtheorem{assumption}{Assumption}
\theoremstyle{remark}
\newtheorem{remark}[theorem]{Remark}

% --- notation ---------------------------------------------------------------------
% One macro per symbol in docs/README_math.md section 1.

\newcommand{\R}{\mathbb{R}}
\newcommand{\Nbr}{\mathcal{N}}                   % open neighborhood (subscript written as \Nbr_i)
\newcommand{\Nbrc}{\bar{\mathcal{N}}}            % closed neighborhood (subscript written as \Nbrc_i)
\newcommand{\ycopy}[2]{y_{#1}^{#2}}             % local copy
\newcommand{\zcon}[1]{z^{#1}}                   % consensus variable
\newcommand{\dual}[2]{\lambda_{#1}^{#2}}        % scaled dual (code `lam`)
\newcommand{\udual}[2]{\nu_{#1}^{#2}}           % unscaled dual = rho * scaled
\newcommand{\Lagr}{\mathcal{L}}
\newcommand{\Lrho}{\Lagr_\rho}
\newcommand{\norm}[1]{\left\lVert #1 \right\rVert}
\newcommand{\prox}{\operatorname{prox}}
\newcommand{\grad}{\nabla}
\newcommand{\diag}{\operatorname{diag}}
\newcommand{\tr}{\operatorname{tr}}
\newcommand{\Eset}{\mathcal{E}}                 % formation-edge set
\newcommand{\Gset}{\mathcal{G}}                 % communication graph
\newcommand{\Fiedler}{\lambda_2(L)}             % algebraic connectivity
\newcommand{\lmax}{\lambda_{\max}}
\newcommand{\lmin}{\lambda_{\min}}
\DeclareMathOperator*{\argmin}{arg\,min}
\DeclareMathOperator*{\argmax}{arg\,max}
. Define macros here only -- never in an individual file,
% or the three documents will drift into inconsistent notation.
%
% Notation single-source-of-truth: docs/README_math.md section 1. Every macro below
% corresponds one-to-one to a row of that table; do not add a symbol there without a
% macro here (or vice versa).

\usepackage[margin=1in]{geometry}
\usepackage{amsmath, amssymb, amsthm}
\usepackage{mathtools}
\usepackage{bm}
\usepackage{algorithm}
\usepackage{algpseudocode}
\usepackage{hyperref}
\usepackage{cleveref}

% --- theorem environments ---------------------------------------------------------
\theoremstyle{plain}
\newtheorem{theorem}{Theorem}
\newtheorem{lemma}[theorem]{Lemma}
\newtheorem{proposition}[theorem]{Proposition}
\newtheorem{corollary}[theorem]{Corollary}
\theoremstyle{definition}
\newtheorem{definition}[theorem]{Definition}
\newtheorem{assumption}{Assumption}
\theoremstyle{remark}
\newtheorem{remark}[theorem]{Remark}

% --- notation ---------------------------------------------------------------------
% One macro per symbol in docs/README_math.md section 1.

\newcommand{\R}{\mathbb{R}}
\newcommand{\Nbr}{\mathcal{N}}                   % open neighborhood (subscript written as \Nbr_i)
\newcommand{\Nbrc}{\bar{\mathcal{N}}}            % closed neighborhood (subscript written as \Nbrc_i)
\newcommand{\ycopy}[2]{y_{#1}^{#2}}             % local copy
\newcommand{\zcon}[1]{z^{#1}}                   % consensus variable
\newcommand{\dual}[2]{\lambda_{#1}^{#2}}        % scaled dual (code `lam`)
\newcommand{\udual}[2]{\nu_{#1}^{#2}}           % unscaled dual = rho * scaled
\newcommand{\Lagr}{\mathcal{L}}
\newcommand{\Lrho}{\Lagr_\rho}
\newcommand{\norm}[1]{\left\lVert #1 \right\rVert}
\newcommand{\prox}{\operatorname{prox}}
\newcommand{\grad}{\nabla}
\newcommand{\diag}{\operatorname{diag}}
\newcommand{\tr}{\operatorname{tr}}
\newcommand{\Eset}{\mathcal{E}}                 % formation-edge set
\newcommand{\Gset}{\mathcal{G}}                 % communication graph
\newcommand{\Fiedler}{\lambda_2(L)}             % algebraic connectivity
\newcommand{\lmax}{\lambda_{\max}}
\newcommand{\lmin}{\lambda_{\min}}
\DeclareMathOperator*{\argmin}{arg\,min}
\DeclareMathOperator*{\argmax}{arg\,max}
. Define macros here only -- never in an individual file,
% or the three documents will drift into inconsistent notation.
%
% Notation single-source-of-truth: docs/README_math.md section 1. Every macro below
% corresponds one-to-one to a row of that table; do not add a symbol there without a
% macro here (or vice versa).

\usepackage[margin=1in]{geometry}
\usepackage{amsmath, amssymb, amsthm}
\usepackage{mathtools}
\usepackage{bm}
\usepackage{algorithm}
\usepackage{algpseudocode}
\usepackage{hyperref}
\usepackage{cleveref}

% --- theorem environments ---------------------------------------------------------
\theoremstyle{plain}
\newtheorem{theorem}{Theorem}
\newtheorem{lemma}[theorem]{Lemma}
\newtheorem{proposition}[theorem]{Proposition}
\newtheorem{corollary}[theorem]{Corollary}
\theoremstyle{definition}
\newtheorem{definition}[theorem]{Definition}
\newtheorem{assumption}{Assumption}
\theoremstyle{remark}
\newtheorem{remark}[theorem]{Remark}

% --- notation ---------------------------------------------------------------------
% One macro per symbol in docs/README_math.md section 1.

\newcommand{\R}{\mathbb{R}}
\newcommand{\Nbr}{\mathcal{N}}                   % open neighborhood (subscript written as \Nbr_i)
\newcommand{\Nbrc}{\bar{\mathcal{N}}}            % closed neighborhood (subscript written as \Nbrc_i)
\newcommand{\ycopy}[2]{y_{#1}^{#2}}             % local copy
\newcommand{\zcon}[1]{z^{#1}}                   % consensus variable
\newcommand{\dual}[2]{\lambda_{#1}^{#2}}        % scaled dual (code `lam`)
\newcommand{\udual}[2]{\nu_{#1}^{#2}}           % unscaled dual = rho * scaled
\newcommand{\Lagr}{\mathcal{L}}
\newcommand{\Lrho}{\Lagr_\rho}
\newcommand{\norm}[1]{\left\lVert #1 \right\rVert}
\newcommand{\prox}{\operatorname{prox}}
\newcommand{\grad}{\nabla}
\newcommand{\diag}{\operatorname{diag}}
\newcommand{\tr}{\operatorname{tr}}
\newcommand{\Eset}{\mathcal{E}}                 % formation-edge set
\newcommand{\Gset}{\mathcal{G}}                 % communication graph
\newcommand{\Fiedler}{\lambda_2(L)}             % algebraic connectivity
\newcommand{\lmax}{\lambda_{\max}}
\newcommand{\lmin}{\lambda_{\min}}
\DeclareMathOperator*{\argmin}{arg\,min}
\DeclareMathOperator*{\argmax}{arg\,max}
